\documentclass{article}
\usepackage{amsmath,amssymb,amsthm,algorithm,algorithmic}
\usepackage{bm}
\usepackage{enumitem}
\usepackage{hyperref}
\usepackage{geometry}
\geometry{margin=1in}
\newtheorem{theorem}{Theorem}
\newtheorem{lemma}{Lemma}
\newtheorem{assumption}{Assumption}
\newcommand{\E}{\mathbb{E}}
\newcommand{\R}{\mathbb{R}}
\newcommand{\grad}{\nabla}
\newcommand{\e}{\mathbf{e}}
\newcommand{\p}{\mathbf{p}}
\newcommand{\w}{\mathbf{w}}
\newcommand{\g}{\mathbf{g}}
\begin{document}

\section*{Randomized Coordinate Descent with Probability Weighting (RCD-PW)}

\section*{Mathematical Formulation}
Consider a convex differentiable function 
\[
f:\R^{d}\rightarrow\R,\qquad f\in C^{1}(\R^{d}),
\]
with coordinate-wise $L_{i}$‑smoothness:

\[
\bigl|\partial_{i}f(\mathbf{x}+ \alpha \mathbf{e}_{i})-\partial_{i}f(\mathbf{x})\bigr|
\le L_{i}\,|\alpha|,\qquad \forall \mathbf{x}\in\R^{d},\;\forall \alpha\in\R,\; i\in[d].
\tag{1}
\]

Let $\p=(p_{1},\dots,p_{d})$ be a probability vector with $p_{i}>0$ and $\sum_{i=1}^{d}p_{i}=1$.
At iteration $k$ we draw an index $i_{k}\in[d]$ independently according to $\p$,
i.e. $\Pr(i_{k}=i)=p_{i}$.  
Define the stochastic coordinate gradient
\[
\g_{k}= \frac{1}{p_{i_{k}}}\,\partial_{i_{k}}f(\w_{k})\,\e_{i_{k}}\in\R^{d},
\tag{2}
\]
where $\e_{i}$ denotes the $i$‑th canonical basis vector.  
The update rule is
\[
\boxed{\;
\w_{k+1}= \w_{k}-\eta\,\g_{k}
      =\w_{k}-\frac{\eta}{p_{i_{k}}}\,\partial_{i_{k}}f(\w_{k})\,\e_{i_{k}}
\;}
\tag{3}
\]
with a constant stepsize $\eta>0$ (later we will restrict $\eta$).

\section*{Pseudocode}
\begin{algorithm}[H]
\caption{Randomized Coordinate Descent with Probability Weighting (RCD‑PW)}
\begin{algorithmic}[1]
\STATE \textbf{Input:} initial point $\w_{0}\in\R^{d}$, stepsize $\eta>0$, probability vector $\p$.
\FOR{$k=0,1,\dots,T-1$}
    \STATE Sample $i_{k}\in[d]$ with $\Pr(i_{k}=i)=p_{i}$.
    \STATE Compute scalar gradient $g_{k}= \partial_{i_{k}}f(\w_{k})$.
    \STATE Update coordinate $i_{k}$:
    \[
        \w_{k+1}= \w_{k}-\frac{\eta}{p_{i_{k}}}\,g_{k}\,\e_{i_{k}}.
    \]
\ENDFOR
\STATE \textbf{Output:} $\w_{T}$ (or the averaged iterate $\bar{\w}_{T}= \frac1T\sum_{k=0}^{T-1}\w_{k}$).
\end{algorithmic}
\end{algorithm}

\section*{Dry Run Example}
We minimise the simple one‑dimensional convex function 
\[
f(w)=(w-3)^{2},\qquad w\in\R .
\]
Here $d=1$, $p_{1}=1$, $L_{1}=2$ (since $f''(w)=2$).  
Take $\eta=0.1$, $w_{0}=0$.

\begin{center}
\begin{tabular}{c|c|c|c|c}
$k$ & $i_{k}$ & $g_{k}= \partial_{i_{k}}f(w_{k})$ & $w_{k+1}=w_{k}-\frac{\eta}{p_{i_{k}}}g_{k}$ & $f(w_{k+1})$\\ \hline
0 & 1 & $2(w_{0}-3)=-6$ & $0-0.1\cdot(-6)=0.6$ & $(0.6-3)^{2}=5.76$\\
1 & 1 & $2(0.6-3)=-4.8$ & $0.6-0.1\cdot(-4.8)=1.08$ & $(1.08-3)^{2}=3.6864$\\
2 & 1 & $2(1.08-3)=-3.84$ & $1.08-0.1\cdot(-3.84)=1.464$ & $(1.464-3)^{2}=2.3660$\\
\end{tabular}
\end{center}
The objective value decreases at each step, illustrating the algorithm.

\newpage
\section*{Assumptions}
\begin{assumption}[Convexity and Differentiability]\label{as:convex}
$f$ is convex and continuously differentiable on $\R^{d}$.
\end{assumption}
\begin{assumption}[Coordinate‑wise $L_{i}$‑smoothness]\label{as:smooth}
For each $i\in[d]$, inequality \eqref{1} holds with a finite constant $L_{i}>0$.
\end{assumption}
\begin{assumption}[Probability vector]\label{as:prob}
$p_{i}>0$ for all $i$ and $\sum_{i=1}^{d}p_{i}=1$.
\end{assumption}
\begin{assumption}[Stepsize]\label{as:stepsize}
The stepsize satisfies $0<\eta\le \min_{i}\frac{p_{i}}{L_{i}}$.
\end{assumption}

\section*{Main Convergence Theorem}
\begin{theorem}[Expected $O(1/k)$ convergence]\label{thm:rate}
Let Assumptions \ref{as:convex}–\ref{as:stepsize} hold.
Define $L_{\max}:=\max_{i}L_{i}$ and $p_{\min}:=\min_{i}p_{i}$.
For the iterates $\{\w_{k}\}$ generated by \eqref{3} we have for all $K\ge1$,
\[
\E\!\left[f(\w_{K})-f(\w^{\star})\right]
\;\le\;
\frac{\|\w_{0}-\w^{\star}\|^{2}}{2\eta K},
\qquad 
\w^{\star}\in\arg\min_{\w}f(\w).
\tag{4}
\]
Consequently, the averaged iterate $\bar{\w}_{K}:=\frac1K\sum_{k=0}^{K-1}\w_{k}$ satisfies
\[
\E\!\left[f(\bar{\w}_{K})-f(\w^{\star})\right]\le \frac{\|\w_{0}-\w^{\star}\|^{2}}{2\eta K}.
\]
\end{theorem}

\section*{Theoretical Properties}
\begin{itemize}[leftmargin=*]
\item \textbf{Unbiasedness.} By construction of $\g_{k}$,
\[
\E[\g_{k}\mid\w_{k}]
= \sum_{i=1}^{d}p_{i}\frac{1}{p_{i}}\partial_{i}f(\w_{k})\e_{i}
= \grad f(\w_{k}).
\]
Thus the stochastic direction is an unbiased estimator of the full gradient.
\item \textbf{Stability w.r.t.\ $p_{i}$.} The stepsize bound \eqref{as:stepsize} scales with $p_{i}$,
ensuring that coordinates sampled rarely (small $p_{i}$) are compensated by a larger
effective step $\eta/p_{i}$, while still guaranteeing monotone expected descent.
\item \textbf{Comparison.} Standard (uniform) coordinate descent corresponds to $p_{i}=1/d$.
Our result reduces to the classical $O(1/k)$ rate with $\eta\le \frac{1}{dL_{\max}}$.
Choosing non‑uniform $p_{i}\propto L_{i}^{-1}$ yields the larger permissible stepsize
$\eta\le \big(\sum_{i}L_{i}^{-1}\big)^{-1}$ and often a tighter bound.
\end{itemize}

\section*{Discussion}
The theorem matches the optimal sublinear rate for first‑order methods on convex smooth functions.
The proof hinges on a careful weighting of expectations by $1/p_{i}$, which eliminates the bias
introduced by non‑uniform sampling. The same technique extends to proximal coordinate updates
and to strongly convex settings, where a linear rate $O((1-\mu\eta)^{k})$ can be obtained.

\newpage
\section*{Preliminaries (Definitions and Lemmas)}
\begin{lemma}[Coordinate smoothness inequality]\label{lem:smooth}
For any $i\in[d]$, any $\w\in\R^{d}$ and any scalar $\alpha\in\R$,
\[
f(\w+\alpha\e_{i})\le f(\w)+\alpha\,\partial_{i}f(\w)+\frac{L_{i}}{2}\alpha^{2}.
\tag{5}
\]
\end{lemma}
\begin{proof}
Apply the definition \eqref{1} and integrate the gradient bound; the standard one‑dimensional
smoothness inequality yields \eqref{5}. ∎
\end{proof}

\begin{lemma}[Unbiased stochastic direction]\label{lem:unbiased}
For the random vector $\g_{k}$ defined in \eqref{2},
\[
\E[\g_{k}\mid\w_{k}]=\grad f(\w_{k}).
\tag{6}
\]
\end{lemma}
\begin{proof}
\[
\E[\g_{k}\mid\w_{k}]
= \sum_{i=1}^{d}p_{i}\frac{1}{p_{i}}\partial_{i}f(\w_{k})\e_{i}
= \sum_{i=1}^{d}\partial_{i}f(\w_{k})\e_{i}
= \grad f(\w_{k}).\quad\text{∎}
\]
\end{proof}

\section*{Main Proof (Step‑by‑Step Derivation)}
We prove Theorem \ref{thm:rate}. Let $\w^{\star}\in\arg\min f$.  
Define the squared distance $D_{k}:=\|\w_{k}-\w^{\star}\|^{2}$.

\noindent\textbf{[Step 1]} Expand $D_{k+1}$ using the update \eqref{3}:
\[
\begin{aligned}
D_{k+1}
&= \bigl\|\w_{k}-\eta\g_{k}-\w^{\star}\bigr\|^{2}\\
&= \|\w_{k}-\w^{\star}\|^{2}
   -2\eta\langle\g_{k},\w_{k}-\w^{\star}\rangle
   +\eta^{2}\|\g_{k}\|^{2}.
\end{aligned}
\tag{7}
\]

\noindent\textbf{[Step 2]} Take conditional expectation w.r.t.\ $i_{k}$ (equivalently w.r.t.\ $\g_{k}$) and use Lemma \ref{lem:unbiased}:
\[
\begin{aligned}
\E\!\bigl[D_{k+1}\mid\w_{k}\bigr]
&= D_{k}
   -2\eta\bigl\langle \E[\g_{k}\mid\w_{k}],\w_{k}-\w^{\star}\bigr\rangle
   +\eta^{2}\E\!\bigl[\|\g_{k}\|^{2}\mid\w_{k}\bigr] \\
&= D_{k}
   -2\eta\langle\grad f(\w_{k}),\w_{k}-\w^{\star}\rangle
   +\eta^{2}\E\!\bigl[\|\g_{k}\|^{2}\mid\w_{k}\bigr].
\end{aligned}
\tag{8}
\]

\noindent\textbf{[Step 3]} Because only coordinate $i_{k}$ is non‑zero,
\[
\|\g_{k}\|^{2}
= \frac{1}{p_{i_{k}}^{2}}\bigl(\partial_{i_{k}}f(\w_{k})\bigr)^{2}.
\]
Hence
\[
\E\!\bigl[\|\g_{k}\|^{2}\mid\w_{k}\bigr]
= \sum_{i=1}^{d}p_{i}\frac{1}{p_{i}^{2}}\bigl(\partial_{i}f(\w_{k})\bigr)^{2}
= \sum_{i=1}^{d}\frac{1}{p_{i}}\bigl(\partial_{i}f(\w_{k})\bigr)^{2}.
\tag{9}
\]

\noindent\textbf{[Step 4]} Use convexity of $f$:
\[
f(\w_{k})-f(\w^{\star})\le \langle\grad f(\w_{k}),\w_{k}-\w^{\star}\rangle .
\tag{10}
\]

\noindent\textbf{[Step 5]} Substitute \eqref{10} into \eqref{8} and rearrange:
\[
\begin{aligned}
\E\!\bigl[D_{k+1}\mid\w_{k}\bigr]
&\le D_{k}
   -2\eta\bigl(f(\w_{k})-f(\w^{\star})\bigr)
   +\eta^{2}\sum_{i=1}^{d}\frac{1}{p_{i}}\bigl(\partial_{i}f(\w_{k})\bigr)^{2}.
\end{aligned}
\tag{11}
\]

\noindent\textbf{[Step 6]} Apply the coordinate‑wise smoothness inequality \eqref{5} with
$\alpha=-\eta/p_{i}$ for each $i$ and sum weighted by $p_{i}$:
\[
\begin{aligned}
f(\w_{k}) 
&\le \sum_{i=1}^{d}p_{i}\Bigl[
      f(\w_{k})-\frac{\eta}{p_{i}}\partial_{i}f(\w_{k})\partial_{i}f(\w_{k})
      +\frac{L_{i}}{2}\Bigl(\frac{\eta}{p_{i}}\Bigr)^{2}
      \bigl(\partial_{i}f(\w_{k})\bigr)^{2}
      \Bigr]   \\
&= f(\w_{k})-\eta\sum_{i=1}^{d}\frac{1}{p_{i}}\bigl(\partial_{i}f(\w_{k})\bigr)^{2}
   +\frac{\eta^{2}}{2}\sum_{i=1}^{d}\frac{L_{i}}{p_{i}^{2}}
      \bigl(\partial_{i}f(\w_{k})\bigr)^{2}.
\end{aligned}
\tag{12}
\]

\noindent\textbf{[Step 7]} Rearranging \eqref{12} yields a bound on the weighted gradient sum:
\[
\sum_{i=1}^{d}\frac{1}{p_{i}}\bigl(\partial_{i}f(\w_{k})\bigr)^{2}
\le \frac{1}{\eta}\bigl(f(\w_{k})-f(\w_{k+1})\bigr)
   +\frac{\eta}{2}\sum_{i=1}^{d}\frac{L_{i}}{p_{i}^{2}}
      \bigl(\partial_{i}f(\w_{k})\bigr)^{2}.
\tag{13}
\]

\noindent\textbf{[Step 8]} Plug \eqref{13} into the RHS of \eqref{11}:
\[
\begin{aligned}
\E\!\bigl[D_{k+1}\mid\w_{k}\bigr]
&\le D_{k}
   -2\eta\bigl(f(\w_{k})-f(\w^{\star})\bigr) \\
&\quad +\eta^{2}\Bigl[
      \frac{1}{\eta}\bigl(f(\w_{k})-f(\w_{k+1})\bigr)
      +\frac{\eta}{2}\sum_{i=1}^{d}\frac{L_{i}}{p_{i}^{2}}
          \bigl(\partial_{i}f(\w_{k})\bigr)^{2}
      \Bigr] \\
&= D_{k}
   -2\eta\bigl(f(\w_{k})-f(\w^{\star})\bigr)
   +\eta\bigl(f(\w_{k})-f(\w_{k+1})\bigr)
   +\frac{\eta^{2}}{2}\sum_{i=1}^{d}\frac{L_{i}}{p_{i}^{2}}
          \bigl(\partial_{i}f(\w_{k})\bigr)^{2}.
\end{aligned}
\tag{14}
\]

\noindent\textbf{[Step 9]} By Assumption \ref{as:stepsize} we have $\eta L_{i}\le p_{i}$,
i.e. $\frac{\eta L_{i}}{p_{i}^{2}}\le\frac{1}{p_{i}}$. Consequently,
\[
\frac{\eta^{2}}{2}\sum_{i=1}^{d}\frac{L_{i}}{p_{i}^{2}}
          \bigl(\partial_{i}f(\w_{k})\bigr)^{2}
\le \frac{\eta}{2}\sum_{i=1}^{d}\frac{1}{p_{i}}
          \bigl(\partial_{i}f(\w_{k})\bigr)^{2}.
\tag{15}
\]

\noindent\textbf{[Step 10]} Using \eqref{13} again to bound the weighted sum in \eqref{15}:
\[
\frac{\eta}{2}\sum_{i=1}^{d}\frac{1}{p_{i}}
          \bigl(\partial_{i}f(\w_{k})\bigr)^{2}
\le \frac{1}{2}\bigl(f(\w_{k})-f(\w_{k+1})\bigr)
   +\frac{\eta^{2}}{4}\sum_{i=1}^{d}\frac{L_{i}}{p_{i}^{2}}
          \bigl(\partial_{i}f(\w_{k})\bigr)^{2}.
\tag{16}
\]

\noindent\textbf{[Step 11]} Substituting \eqref{16} into \eqref{14} and discarding the non‑negative
term $\frac{\eta^{2}}{4}\sum_{i}\frac{L_{i}}{p_{i}^{2}}(\partial_{i}f)^{2}$ yields the clean recursion:
\[
\E\!\bigl[D_{k+1}\mid\w_{k}\bigr]
\le D_{k}
   -\eta\bigl(f(\w_{k})-f(\w^{\star})\bigr)
   -\frac{\eta}{2}\bigl(f(\w_{k})-f(\w_{k+1})\bigr).
\tag{17}
\]

\noindent\textbf{[Step 12]} Take total expectation and sum \eqref{17} over $k=0,\dots,K-1$:
\[
\begin{aligned}
\E[D_{K}]
&\le D_{0}
   -\eta\sum_{k=0}^{K-1}\E\!\bigl[f(\w_{k})-f(\w^{\star})\bigr]
   -\frac{\eta}{2}\sum_{k=0}^{K-1}\E\!\bigl[f(\w_{k})-f(\w_{k+1})\bigr] \\
&= D_{0}
   -\eta\sum_{k=0}^{K-1}\E\!\bigl[f(\w_{k})-f(\w^{\star})\bigr]
   -\frac{\eta}{2}\Bigl(\E[f(\w_{0})]-\E[f(\w_{K})]\Bigr).
\end{aligned}
\tag{18}
\]

\noindent\textbf{[Step 13]} Drop the non‑negative term $-\frac{\eta}{2}\bigl(\E[f(\w_{0})]-\E[f(\w_{K})]\bigr)$
and use $D_{K}\ge0$ to obtain
\[
\eta\sum_{k=0}^{K-1}\E\!\bigl[f(\w_{k})-f(\w^{\star})\bigr]\le D_{0}.
\tag{19}
\]

\noindent\textbf{[Step 14]} Divide by $\eta K$ and use convexity of $f$ together with Jensen’s inequality
to bound the average iterate:
\[
\E\!\bigl[f(\bar{\w}_{K})-f(\w^{\star})\bigr]
\le \frac{1}{K}\sum_{k=0}^{K-1}\E\!\bigl[f(\w_{k})-f(\w^{\star})\bigr]
\le \frac{D_{0}}{2\eta K}.
\tag{20}
\]
Since $D_{0}=\|\w_{0}-\w^{\star}\|^{2}$, inequality \eqref{20} is exactly the claim \eqref{4}.

Thus Theorem \ref{thm:rate} is proved. ∎

\section*{Rate Derivation (With Constants)}
From the steps above we see that the constant in \eqref{4} is $1/(2\eta)$.  
Choosing the largest admissible stepsize $\eta^{\star}= \min_{i}\frac{p_{i}}{L_{i}}$ gives
\[
\E\!\bigl[f(\bar{\w}_{K})-f(\w^{\star})\bigr]
\le
\frac{L_{\max}}{2p_{\min}}\,
\frac{\|\w_{0}-\w^{\star}\|^{2}}{K}.
\]
If $p_{i}\propto 1/L_{i}$ (optimal importance sampling) then $p_{\min}=1/\sum_{j}L_{j}$ and the bound becomes
\[
\E\!\bigl[f(\bar{\w}_{K})-f(\w^{\star})\bigr]
\le
\frac{\bigl(\sum_{j=1}^{d}L_{j}\bigr)}{2}\,
\frac{\|\w_{0}-\w^{\star}\|^{2}}{K},
\]
which can be substantially smaller than the uniform‑sampling bound that involves $dL_{\max}$.

\section*{Special Cases}
\begin{itemize}[leftmargin=*]
\item \textbf{Uniform sampling.} $p_{i}=1/d$ yields $\eta\le \frac{1}{dL_{\max}}$ and recovers the classic $O(dL_{\max}/K)$ rate.
\item \textbf{Importance sampling.} Setting $p_{i}=L_{i}^{-1}\big/\sum_{j}L_{j}^{-1}$ maximises the admissible stepsize and leads to the bound with $\sum_{j}L_{j}$ instead of $dL_{\max}$.
\item \textbf{Strong convexity.} If $f$ is $\mu$‑strongly convex, the same proof with the stronger inequality 
$f(\w_{k})-f(\w^{\star})\ge \frac{\mu}{2}\|\w_{k}-\w^{\star}\|^{2}$ yields a linear rate
$\E[f(\w_{k})-f(\w^{\star})]\le (1-\mu\eta)^{k}\bigl(f(\w_{0})-f(\w^{\star})\bigr)$.
\end{itemize}

\end{document}